% Generated draft: requires mathematical and source review.
\documentclass[11pt]{article}
\usepackage[margin=25mm]{geometry}
\usepackage{fontspec}
\setmainfont{Latin Modern Roman}
\usepackage{amsmath,amssymb,mathtools,unicode-math}
\setmathfont{Latin Modern Math}
\usepackage{enumitem,needspace,xurl,hyperref}
\setlength{\parindent}{0pt}
\setlength{\parskip}{5pt}
\setlength{\emergencystretch}{3em}
\newcommand{\N}{\mathbb{N}}
\newcommand{\Z}{\mathbb{Z}}
\newcommand{\Q}{\mathbb{Q}}
\newcommand{\R}{\mathbb{R}}
\newcommand{\C}{\mathbb{C}}
\newcommand{\F}{\mathbb{F}}
\newcommand{\A}{\mathbb{A}}
\newcommand{\PP}{\mathbb{P}}
\newcommand{\E}{\mathbb{E}}
\newcommand{\eps}{\varepsilon}
\newcommand{\dd}{\,\mathrm{d}}
\begin{document}
\section*{UnsolvedMath ID: 30005508}
\textbf{Source problem code:} catalogue entry
\subsection{Definitions and mathematical statement}
Let \(G\) be a finite group. Use a splitting \(2\)-modular system for the subgroups involved, so that projective indecomposable characters are ordinary characters of lifts of indecomposable projective modules in characteristic \(2\). A \(2\)-block \(B\) of \(G\) is called real if complex conjugation preserves its set \(\operatorname{Irr}(B)\) of ordinary irreducible characters; it is nonprincipal if it does not contain the trivial character of \(G\).

Let \(B\) be a real nonprincipal \(2\)-block with a fixed defect pair \((D,E)\), where \(D\) is its defect group and \(E\) its extended defect group. The term ``defect pair'' is used in the real-block-theoretic sense of the supplied source. The record does not supply its defining construction or conventions; these must be checked against the cited report, and the defect-pair hypotheses below must be retained in full.

For a subgroup \(H\leq G\) and \(g\in G\), write \(C_H(g)=\{h\in H:hg=gh\}\). Let \((x,b)\) be a \(B\)-subsection: \(x\) is an element of \(2\)-power order and \(b\) is a block of \(C_G(x)\) associated with \(B\) by the Brauer map at \(\langle x\rangle\). More explicitly, over a splitting field \(k\) of characteristic \(2\), its block idempotent satisfies
\[
 e_b\operatorname{Br}_{\langle x\rangle}(e_B)\ne 0,
\]
where \(\operatorname{Br}_{\langle x\rangle}:(kG)^{\langle x\rangle}\longrightarrow kC_G(x)\) is the Brauer homomorphism. Assume that \(b\), as a block of \(C_G(x)\), has defect pair exactly \((C_D(x),C_E(x))\), and that \(b\) has a unique projective indecomposable character, denoted by \(\Phi\).

For class functions \(\alpha,\beta\) on a finite group \(H\), put
\[
 [\alpha,\beta]_H=\frac{1}{|H|}\sum_{h\in H}\alpha(h)\overline{\beta(h)}.
\]
Write \(1_H\) for the trivial character of \(H\), \(\operatorname{Res}^{K}_{H}\) for restriction from \(K\) to \(H\), and
\[
 y^{C_G(x)}=\{h^{-1}yh:h\in C_G(x)\}.
\]
If \(y^2=x\), then \(y\in C_G(x)\) and \(C_G(y)\leq C_G(x)\), so both the conjugacy class and the restriction below are defined.

\textbf{Statement (local Conjecture 4 in the supplied excerpt).} For every finite group \(G\) and every choice of the preceding data satisfying all the stated hypotheses, every \(y\in G\) with \(y^2=x\) satisfies
\[
 \left[\operatorname{Res}^{C_G(x)}_{C_G(y)}\Phi,1_{C_G(y)}\right]_{C_G(y)}
 =\left|y^{C_G(x)}\cap(E\setminus D)\right|.
\]
The intersection is taken with the set difference \(E\setminus D\). If \(x\) has no square root in \(G\), the assertion for that subsection is vacuous.

\textit{Supplied context and qualifications.} The excerpt states that this local conjecture implies its Conjecture 2, whose statement is absent from the record. Its Theorem 3 asserts that Conjecture 2 holds when \(D\) is abelian or dihedral, or when \(G\) is solvable or quasisimple. Those assertions concern Conjecture 2 and do not by themselves establish the local conjecture above. The excerpt also gives the following consequence of the local conjecture, with \(\operatorname{IBr}(b)=\{\varphi\}\):
\[
 \sum_{\chi\in\operatorname{Irr}(B)}\epsilon(\chi)d^x_{\chi\varphi}
 =\left|\{y\in E\setminus D:y^2=x\}\right|,
\]
where \(\operatorname{IBr}(b)\) denotes the irreducible Brauer characters of \(b\), \(d^x_{\chi\varphi}\) are generalized decomposition numbers, and \(\epsilon(\chi)=|G|^{-1}\sum_{g\in G}\chi(g^2)\) is the Frobenius--Schur indicator. This consequence is supplied context, not a separate requested conjecture. The subsequent Theorem 5 concerns real blocks with dihedral defect groups of order at least \(8\); its classification table is incompletely transcribed in the record and is not used to alter the present statement.
\subsection{Short English statement}
For a real nonprincipal \(2\)-block and a subsection satisfying the specified defect-pair and uniqueness hypotheses, the conjecture identifies the multiplicity of the trivial character in a restricted projective indecomposable character with the number of conjugates of a chosen square root lying in \(E\setminus D\). The record labels the problem open; this is supplied dated curation, not an independently verified status assessment.
\subsection{Sources}
\begin{itemize}
\item Supplied local UnsolvedMath registry record, permanent ID \(30005508\), title ``Projective Characters and Square Roots in Real Blocks,'' exact source problem code \texttt{catalogue entry}, in the Oberwolfach Reports Open Problems collection. The record gives proposed year \(2023\) and does not identify a proposer.
\item \emph{Representations of Finite Groups} (2023), Oberwolfach Reports, report 19/2023, pp.\ 1031--1084, DOI 10.4171/owr/2023/19: \url{https://doi.org/10.4171/owr/2023/19}. The supplied excerpt identifies the target as local Conjecture 4; no narrower page reference is supplied.
\item Supplied literature source: Sambale, \emph{Real characters in nilpotent blocks}, \url{https://doi.org/10.1007/s10013-023-00623-5}.
\item Supplied literature source: \emph{Oberwolfach Report 19/2023}, \url{https://doi.org/10.4171/owr/2023/19}.
\item Literature assessment dated exactly \texttt{2026-08-22}: the supplied curation labels the problem open, describes local Conjecture 4 as a sufficient route to the general indicator formula, and reports that the report contains no proof and that the literature does not establish the conjecture in general. These are supplied dated curation claims, not independently verified findings. The background explicitly requests independent verification of time-sensitive status and cited claims.
\item The registry labels the statement \texttt{corrected\_verified} and says it was reconstructed and checked against the official report. This records the supplied provenance claim; this draft does not independently verify that comparison or assert mathematical correctness.
\end{itemize}
\end{document}
